All stochastic orders below are uniform over \(P\in\mathfrak P_T\) and, when a contrast is present, over \(\|c\|_2=1\). By the bounded-score and uniform-overlap conditions in \(\textnormal{def:finite-attractor-class}\),
\[
 \|\xi_t\|_2\le M_\xi:=\frac{B}{\varepsilon}
 \quad\text{almost surely}. \tag{1}
\]
Define
\[
 D_t=\xi_t\xi_t^\top-E_P[\xi_t\xi_t^\top\mid\mathcal F_{t-1}].
\]
Every entry of \(D_t\) is a martingale difference. Moreover,
\[
 \|D_t\|_{\mathrm F}
 \le \|\xi_t\xi_t^\top\|_{\mathrm F}
   +E_P[\|\xi_t\xi_t^\top\|_{\mathrm F}\mid\mathcal F_{t-1}]
 \le 2M_\xi^2. \tag{2}
\]
Martingale orthogonality, summed over the fixed \(d^2\) entries, therefore gives
\[
 E_P\left\|T^{-1}\sum_{t=1}^TD_t\right\|_{\mathrm F}^2
 =T^{-2}\sum_{t=1}^TE_P\|D_t\|_{\mathrm F}^2
 \le \frac{4M_\xi^4}{T}. \tag{3}
\]
Since the operator norm is bounded by the Frobenius norm, Chebyshev's inequality applied to (3) proves
\[
 T^{-1}([S]_T-\langle S\rangle_T)=O_P(T^{-1/2}). \tag{4}
\]

The exact-root localization and fallback conclusions of \(\textnormal{thm:basin-stable-ipwz}\) imply
\[
 \|\widehat\theta_T-\theta^*\|_2=O_P(T^{-1/2}),
 \qquad
 \partial_\theta G_T(\theta^*)-H=O_P(T^{-1/2}). \tag{5}
\]
For every \(t\), the derivative and overlap bounds give the pathwise inequality
\[
 \|\widehat\xi_t-\xi_t\|_2
 \le \frac{D_1}{\varepsilon}\|\widehat\theta_T-\theta^*\|_2. \tag{6}
\]
Both scores have norm at most \(M_\xi\), so
\[
 \|\widehat\xi_t\widehat\xi_t^\top-
       \xi_t\xi_t^\top\|_{\mathrm{op}}
 \le (\|\widehat\xi_t\|_2+\|\xi_t\|_2)
       \|\widehat\xi_t-\xi_t\|_2
 \le \frac{2BD_1}{\varepsilon^2}
       \|\widehat\theta_T-\theta^*\|_2. \tag{7}
\]
Averaging (7), using (4)--(5), and using the predictable-QV conclusion of \(\textnormal{thm:basin-stable-ipwz}\),
\[
 \begin{aligned}
 \|\widehat\Omega_T-\Omega_J\|_{\mathrm{op}}
 &\le \left\|T^{-1}\sum_{t=1}^T
   (\widehat\xi_t\widehat\xi_t^\top-
    \xi_t\xi_t^\top)\right\|_{\mathrm{op}}\\
 &\quad+|T^{-1}([S]_T-\langle S\rangle_T)\|_{\mathrm{op}}
 +\|T^{-1}\langle S\rangle_T-\Omega_J\|_{\mathrm{op}}\\
 &=O_P(T^{-1/2}+r_T+\delta_T)=O_P(e_T). \tag{8}
 \end{aligned}
\]
Here the expectation bound for the last term in the cited theorem implies the displayed stochastic order by Markov's inequality. On the closed \(\rho\)-ball, the curvature and overlap bounds similarly yield
\[
 \begin{aligned}
 \|\widehat H_T-H\|_{\mathrm{op}}
 &\le \|\partial_\theta G_T(\theta^*)-H\|_{\mathrm{op}}
 +\frac{D_2}{\varepsilon}\|\widehat\theta_T-\theta^*\|_2\\
 &=O_P(T^{-1/2})=O_P(e_T). \tag{9}
 \end{aligned}
\]
Equations (5)--(9) are unaffected by the fallback branch, because its probability tends to zero uniformly and the displayed orders are uniform-in-probability statements.

The singular-value perturbation inequality, (9), and \(\sigma_{\min}(H)\ge h_-\) show that
\[
 P\{\sigma_{\min}(\widehat H_T)<h_-/2\}\longrightarrow0 \tag{10}
\]
uniformly. Likewise, \(\Omega_J\succeq\omega_-I_d\), (8), and the eigenvalue perturbation inequality give
\[
 P\{\lambda_{\min}(\widehat\Omega_T)<\omega_-/2\}\longrightarrow0. \tag{11}
\]
Thus (10)--(11) and the union bound prove
\[
 \sup_{P\in\mathfrak P_T}P(\mathcal E_T^c)\longrightarrow0. \tag{12}
\]
On \(\mathcal E_T\), the inverse is evaluated legally and
\[
 \widehat H_T^{-1}-H^{-1}
 =\widehat H_T^{-1}(H-\widehat H_T)H^{-1},
 \qquad
 \|\widehat H_T^{-1}\|_{\mathrm{op}}\le\frac2{h_-}. \tag{13}
\]
The finite-mixture theorem also gives
\[
 \|H^{-1}\|_{\mathrm{op}}\le h_-^{-1},
 \qquad
 \|\Omega_J\|_{\mathrm{op}}\le M_\Omega:=\frac{KB^2}{\varepsilon}. \tag{14}
\]
Insert and subtract the three sandwich factors one at a time. Equations (8), (9), (13), and (14) then give, on \(\mathcal E_T\),
\[
 \widehat H_T^{-1}\widehat\Omega_T\widehat H_T^{-\top}
 -H^{-1}\Omega_JH^{-\top}=O_P(e_T). \tag{15}
\]
On \(\mathcal E_T^c\), \(\widehat V_T=I_d\), whereas (14) bounds \(V_J\) uniformly. Since (12) makes this branch uniformly negligible, (15) and \(\textnormal{def:basin-covariance}\) prove
\[
 \widehat V_T-V_J=O_P(e_T). \tag{16}
\]

It remains to derive, rather than assume, the pivotal limit. Fix a unit vector \(c\), and define the deterministic, law-specific vector and scalar martingale differences
\[
 a=H^{-\top}c,
 \qquad U_t=a^\top\xi_t,
 \qquad A_T^2=\sum_{t=1}^TE_P[U_t^2\mid\mathcal F_{t-1}],
 \qquad Q_T=\sum_{t=1}^TU_t^2. \tag{17}
\]
The upper Jacobian bound in \(\textnormal{thm:basin-stable-ipwz}\) and the lower singular-value assumption imply
\[
 \frac1{KD_1}\le\|a\|_2\le\frac1{h_-},
 \qquad |U_t|\le L:=\frac{B}{\varepsilon h_-}. \tag{18}
\]
By (17), the predictable-QV approximation, and regime nondegeneracy,
\[
 \frac{A_T^2}{T}
 =a^\top(T^{-1}\langle S\rangle_T)a
 =a^\top\Omega_Ja+o_P(1),
 \qquad
 a^\top\Omega_Ja\ge
 v_0:=\frac{\omega_-}{(KD_1)^2}>0. \tag{19}
\]
Hence \(A_T^2\to\infty\) in probability uniformly. Equation (4) also yields
\[
 \frac{Q_T-A_T^2}{T}
 =a^\top T^{-1}([S]_T-\langle S\rangle_T)a
 =O_P(T^{-1/2}), \tag{20}
\]
so (19)--(20) imply \(Q_T/A_T^2\to_P1\) and \(P(Q_T>0)\to1\) uniformly. With \(\delta=2\), (18)--(19) imply
\[
 A_T^{-4}\sum_{t=1}^TE_P|U_t|^4
 \le A_T^{-4}TL^4\longrightarrow_P0. \tag{21}
\]
Thus every triangular array obtained from an arbitrary sequence \(P_T\in\mathfrak P_T\) and unit contrasts \(c_T\) satisfies all hypotheses of \(\textnormal{lem:saco-self-normalized-martingale-clt}\). It follows that
\[
 \frac{a^\top S_T}{\sqrt{a^\top[S]_Ta}}
 =\frac{\sum_{t=1}^TU_t}{\sqrt{Q_T}}
 \Rightarrow N(0,1). \tag{22}
\]
This convergence is uniform in \(P\) and unit \(c\). Indeed, failure of uniformity would produce a sequence \((P_T,c_T)\) and a positive lower bound on the Kolmogorov distance, contradicting (22) for that triangular array. To pass from weak convergence to Kolmogorov convergence, partition the real line at finitely many continuity points of \(\Phi\), make every intervening \(\Phi\)-increment arbitrarily small, and use monotonicity of distribution functions on each interval; this proves the required uniform cdf bound without an additional theorem.

By \(\textnormal{thm:basin-stable-ipwz}\),
\[
 \sqrt T\,c^\top(\widehat\theta_T-\theta^*)
 =-T^{-1/2}a^\top S_T+O_P(e_T). \tag{23}
\]
Furthermore, (4), (8), (16), and \(c^\top V_Jc=a^\top\Omega_Ja\) imply
\[
 c^\top\widehat V_Tc-T^{-1}a^\top[S]_Ta=O_P(e_T). \tag{24}
\]
By (19)--(20), both quantities compared in (24) are at least \(v_0/2\) with probability tending uniformly to one. Division and square roots are therefore legitimate on an event whose probability tends uniformly to one. Since \(e_T\to0\), equations (22)--(24) and the elementary identity
\[
 |\sqrt x-\sqrt y|=\frac{|x-y|}{\sqrt x+\sqrt y}
\]
give
\[
 T_T^{\mathrm{rqv}}(c)
 =-\frac{a^\top S_T}{\sqrt{a^\top[S]_Ta}}+o_P(1) \tag{25}
\]
uniformly. A direct smoothing argument applied to (22) and (25), together with uniform continuity of \(\Phi\), proves
\[
 \sup_{P\in\mathfrak P_T}\sup_{\|c\|_2=1}\sup_{z\in\mathbb R}
 \left|P\{T_T^{\mathrm{rqv}}(c)\le z\}-\Phi(z)\right|
 \longrightarrow0. \tag{26}
\]

For completeness, (26) supplies the Wald conclusion required for the causal contrast, not merely a distributional label. For every fixed \(\alpha\in(0,1)\), let \(z_{1-\alpha/2}=\Phi^{-1}(1-\alpha/2)\) and set
\[
 I_T(c,\alpha)=
 \left[\widehat\tau_T-z_{1-\alpha/2}
       \sqrt{\frac{c^\top\widehat V_Tc}{T}},
       \widehat\tau_T+z_{1-\alpha/2}
       \sqrt{\frac{c^\top\widehat V_Tc}{T}}\right].
\]
Because \(\tau=c^\top\theta^*\), membership \(\tau\in I_T(c,\alpha)\) is equivalent to \(|T_T^{\mathrm{rqv}}(c)|\le z_{1-\alpha/2}\). Hence (26) gives
\[
 \sup_{P\in\mathfrak P_T}\sup_{\|c\|_2=1}
 \left|P\{\tau\in I_T(c,\alpha)\}-(1-\alpha)\right|
 \longrightarrow0. \tag{27}
\]
No step above uses pointwise convergence \(\eta_t\to\bar\eta_J\). The regime label enters only through the covariance atom \(\Omega_J\), and every approximation to that atom is supplied by the uniform Cesaro bound already proved in \(\textnormal{thm:basin-stable-ipwz}\). Thus the proof remains valid under the broadened finite-Cesaro-regime condition.